\section{Models used in this study}

To address the challenges of personalized fashion recommendation, we systematically evaluate and compare a series of graph neural collaborative filtering architectures. These architectures range from a pure interaction-based collaborative filtering baseline to advanced models that integrate visual and textual modalities via dual-graph propagation, self-supervised bootstrap representation learning, and decoupled content propagation over frozen affinity structures. 

Specifically, this study builds on LightGCN \parencite{he2020lightgcnsimplifyingpoweringgraph}, which processes only user-item interaction signals, as the shared collaborative-filtering backbone, and benchmarks three multimodal graph architectures that extend it: CombiGCN \parencite{10.1007/978-981-97-0669-3_11}, a dual-graph GCN that aggregates signals from both interaction and content-based item similarity graphs; BM3 \parencite{10.1145/3543507.3583251}, a self-supervised model utilizing bootstrap contrastive learning without negative pairs; and FREEDOM \parencite{Zhou_2023}, which decouples collaborative and content propagation using a frozen item-item kNN graph and aligns them via InfoNCE loss. By comparing these architectures under uniform multimodal embedding inputs and fusion strategies, we analyze the impact of different propagation topologies and optimization objectives on fashion recommendation performance.

\subsection{LightGCN}

As established in Section 2.2, LightGCN \parencite{he2020lightgcnsimplifyingpoweringgraph} provides the collaborative-filtering foundation for this study. It operates exclusively on the bipartite user-item interaction graph $\mathcal{G}_{ui} = (\mathcal{U} \cup \mathcal{I}, \mathcal{E})$ and simplifies standard Graph Convolutional Networks (GNNs) by discarding feature transformation matrices and non-linear activations. In this study, LightGCN is not benchmarked as a standalone configuration; rather, it serves as the collaborative-filtering backbone shared by all three evaluated architectures (CombiGCN, BM3, and FREEDOM), each of which extends LightGCN propagation with its own multimodal mechanism.

For the collaborative filtering branch in all compared architectures, LightGCN performs neighborhood aggregation to learn structural representations. The propagation rule, layer combination, and Bayesian Personalized Ranking (BPR) loss optimization follow the mathematical formulations defined in Equations (\ref{eq:prop_user}), (\ref{eq:prop_item}), (\ref{eq:layer_comb}), and (\ref{eq:bpr_loss}) respectively. This backbone isolates the structural collaborative signal upon which the multimodal mechanisms of the evaluated models build, without itself leveraging visual aesthetics or textual metadata.

\subsection{CombiGCN}

CombiGCN \parencite{10.1007/978-981-97-0669-3_11} originally extends the single bipartite graph propagation of LightGCN by introducing an auxiliary user-user similarity graph constructed from user historical interactions. In this study, to enable multimodal fashion recommendation, we adapt the dual-graph GCN propagation of CombiGCN to operate over an auxiliary item-item similarity graph built from visual and textual content features, formulating recommendation as a dual-graph message-passing process over user-item interactions and item-item similarity links. The primary technical motivation of this adaptation is to address user-item interaction sparsity by propagating collaborative preferences along semantic and visual similarity links between items. The overall workflow and architecture of CombiGCN are illustrated in Figure~\ref{fig:combigcn}.

\begin{figure}[H]
	\centering
	\includegraphics[width=0.9\textwidth]{figs/combigcn}
	\caption{Architecture and workflow of CombiGCN, combining the user-item collaborative filtering (CF) branch with the item-item similarity branch via dual-graph GCN layers.}
	\label{fig:combigcn}
\end{figure}

The workflow of CombiGCN is executed in the following steps:
\begin{enumerate}
    \item \textbf{Precomputed Similarity Graph}: Prior to training, an item-item similarity matrix $\mathbf{W} \in \mathbb{R}^{|\mathcal{I}| \times |\mathcal{I}|}$ is constructed at the data layer. Depending on the configuration of \texttt{sim\_type} (\texttt{img\_only}, \texttt{tfidf}, or \texttt{multimodal}), cosine similarities are computed between the item features. To eliminate noisy connections and enforce sparsity, a threshold filter is applied such that $W_{ij} = 0$ if the cosine similarity is below $0.5$. The matrix is then symmetrically normalized as:
    \begin{equation}
        \mathbf{S} = \mathbf{D}_s^{-1/2} \mathbf{W} \mathbf{D}_s^{-1/2},
    \end{equation}
    where $\mathbf{D}_s$ is the diagonal degree matrix of $\mathbf{W}$.
    \item \textbf{Dual-Graph Propagation}: In each GCN layer $k+1$, user representations are updated solely by aggregating observed item embeddings from the user-item interaction graph:
    \begin{equation}
        \mathbf{e}_u^{(k+1)} = \sum_{i \in \mathcal{N}_u} \frac{1}{\sqrt{|\mathcal{N}_u| |\mathcal{N}_i|}} \mathbf{e}_i^{(k)}.
    \end{equation}
    Conversely, item representations are updated by fusing signals from both the user-item interaction graph (CF branch) and the item-item similarity graph (similarity branch):
    \begin{equation}
        \mathbf{e}_i^{(k+1)} = \mathbf{e}_{i, \text{CF}}^{(k+1)} + \mathbf{e}_{i, \text{Sim}}^{(k+1)},
    \end{equation}
    where the collaborative signal is propagated from user nodes:
    \begin{equation}
        \mathbf{e}_{i, \text{CF}}^{(k+1)} = \sum_{u \in \mathcal{N}_i} \frac{1}{\sqrt{|\mathcal{N}_i| |\mathcal{N}_u|}} \mathbf{e}_u^{(k)},
    \end{equation}
    and the content-similarity signal is propagated from neighboring items:
    \begin{equation}
        \mathbf{e}_{i, \text{Sim}}^{(k+1)} = \sum_{j \in \mathcal{S}_i} S_{ij} \mathbf{e}_j^{(k)},
    \end{equation}
    where $\mathcal{S}_i$ denotes the set of similar items connected to item $i$ in the similarity graph.
    \item \textbf{Layer Aggregation and Optimization}: After propagating through a stack of $N$ CombiGCN layers, the layer outputs are combined via mean pooling to generate the final user and item embeddings. The model parameters (initial ID embeddings $\mathbf{E}^{(0)}$) are optimized using the standard pairwise BPR loss defined in Equation (\ref{eq:bpr_loss}).
\end{enumerate}

By performing fusion at the graph propagation level, CombiGCN leverages precomputed similarity graphs to guide the collaborative filtering process, enabling smoother embeddings for items with shared multimodal features.

\subsection{BM3}

BM3 \parencite{10.1145/3543507.3583251} is a self-supervised multimodal recommendation framework designed to alleviate the computational cost of auxiliary GCN graphs and the noise introduced by negative sampling. Instead of performing graph convolutions over item-item similarity matrices, BM3 operates a single bipartite LightGCN propagation for collaborative filtering and aligns the resulting representations with raw multimodal features inside the network using a bootstrap contrastive learning objective. The detailed architecture and data flow of BM3 are depicted in Figure~\ref{fig:bm3}.

\begin{figure}[H]
	\centering
	\includegraphics[width=0.95\textwidth]{figs/bm3}
	\caption{Architecture and workflow of BM3, featuring the collaborative filtering (CF) branch, linear modal projectors, and the self-supervised bootstrap contrastive learning process with an EMA target encoder.}
	\label{fig:bm3}
\end{figure}

The workflow of BM3 consists of the following components:
\begin{enumerate}
    \item \textbf{ID Embedding and CF Branch}: The model initializes learnable user ID embeddings $\mathbf{h}_u^0 \in \mathbb{R}^{N_u \times d}$ and item ID embeddings $\mathbf{h}_i^0 \in \mathbb{R}^{N_i \times d}$ (with $d = 512$, $N_u = 553$, and $N_i = 2194$ in this study). The initial ego embedding matrix is constructed by concatenating these matrices ($\mathbf{E}^{(0)} \in \mathbb{R}^{(N_u + N_i) \times d}$, size $2747 \times 512$). It is propagated through a 4-layer GCN over the user-item interaction graph (augmented with edge dropout) to yield the collaborative representations $\mathbf{e}_{u, \text{CF}} \in \mathbb{R}^{553 \times 512}$ and $\mathbf{e}_{i, \text{CF}} \in \mathbb{R}^{2194 \times 512}$.
    \item \textbf{Modal Projection and Fusion}: Raw multimodal features, including item image embeddings $\mathbf{x}_{vis,i} \in \mathbb{R}^{2194 \times d_v}$ and text embeddings $\mathbf{x}_{txt,i} \in \mathbb{R}^{2194 \times d_t}$, are processed through linear projection layers to align their dimensionality with the CF space. Here $d_v$ denotes the visual feature dimension ($512$ for CLIP, and $768$ for MobileNetV2 after PCA reduction from its native $1280$-dimensional output), and $d_t$ denotes the TF-IDF vocabulary dimension:
    \begin{equation}
        \mathbf{h}_{v,i} = \mathbf{W}_v \mathbf{x}_{vis,i} + \mathbf{b}_v, \quad \mathbf{h}_{t,i} = \mathbf{W}_t \mathbf{x}_{txt,i} + \mathbf{b}_t,
    \end{equation}
    where $\mathbf{W}_v \in \mathbb{R}^{512 \times d_v}$ and $\mathbf{W}_t \in \mathbb{R}^{512 \times d_t}$ are trainable matrices. The projected features are fused according to the \texttt{sim\_type} parameter to obtain the modal representation $\mathbf{e}_{i, \text{modal}} \in \mathbb{R}^{2194 \times 512}$ (e.g., $\mathbf{e}_{i, \text{modal}} = (\mathbf{h}_{v,i} + \mathbf{h}_{t,i})/2$ for the \texttt{multimodal} configuration, or via an attention layer for \texttt{mm\_attention}).
    \item \textbf{Bootstrap Contrastive Learning}: To align the collaborative and content-based item spaces without negative samples, BM3 utilizes a bootstrap alignment objective. An online predictor head $q(\cdot)$ (consisting of a two-layer MLP with ReLU activation) is applied to the online representations to break symmetry and prevent representation collapse. A target item encoder maintains target embeddings $\mathbf{h}_{i\_target}^0 \in \mathbb{R}^{2194 \times 512}$ and is updated using Exponential Moving Average (EMA) from the online encoder:
    \begin{equation}
        \mathbf{\Theta}_{target} \leftarrow m \mathbf{\Theta}_{target} + (1 - m) \mathbf{\Theta}_{online},
    \end{equation}
    where $m = 0.995$ is the momentum parameter, and gradients are not backpropagated through the target branch. The bootstrap loss evaluates the cosine similarity between the online and target representations:
    \begin{equation}
        \mathcal{L}_{boot}(\mathbf{p}, \mathbf{z}) = 2 - 2 \frac{\mathbf{p}^T \mathbf{z}}{\|\mathbf{p}\| \|\mathbf{z}\|}.
    \end{equation}
    The self-supervised contrastive objective is formulated as:
    \begin{equation}
        \mathcal{L}_{CL} = \frac{1}{2} \left[ \mathcal{L}_{boot}(q(\mathbf{e}_{i, \text{CF}}), sg(\mathbf{e}_{i, \text{modal}})) + \mathcal{L}_{boot}(q(\mathbf{e}_{i, \text{modal}}), sg(\mathbf{e}_{i, \text{target}})) \right],
    \end{equation}
    where $sg(\cdot)$ is the stop-gradient operator.
    \item \textbf{Representation Fusion and Objective}: The final item embedding for user preference scoring is computed as the sum of the collaborative and modal representations:
    \begin{equation}
        \mathbf{h}_i = \mathbf{e}_{i, \text{CF}} + \mathbf{e}_{i, \text{modal}}.
    \end{equation}
    The model is optimized using a joint training objective:
    \begin{equation}
        \mathcal{L}_{total} = \mathcal{L}_{BPR} + \lambda_1 \mathcal{L}_{reg} + \lambda_2 \mathcal{L}_{CL},
    \end{equation}
    where $\mathcal{L}_{BPR}$ is the Bayesian Personalized Ranking loss computed on the fused embeddings ($\mathbf{h}_u^T \mathbf{h}_i$), $\mathcal{L}_{reg}$ is $L_2$ regularization, and $\lambda_2$ is the contrastive loss weight (set to $0.2$ in our configuration).
\end{enumerate}

\subsection{FREEDOM}

FREEDOM \parencite{Zhou_2023} proposes a decoupled graph structure learning framework that separates collaborative filtering propagation and content propagation. Rather than updating item-item relationships dynamically or aligning representations via bootstrap, FREEDOM constructs a static kNN affinity graph at initialization to denoise structural learning. It aligns the resulting decoupled views using a classical InfoNCE contrastive loss that actively contrasts positive item pairs against in-batch negative pairs. The architecture and workflow of FREEDOM are shown in Figure~\ref{fig:freedom}.

\begin{figure}[H]
	\centering
	\includegraphics[width=0.95\textwidth]{figs/freedom}
	\caption{Architecture and workflow of FREEDOM, showing the decoupled interactive view (CF branch) and semantic view (content branch over a frozen kNN graph) aligned via InfoNCE contrastive loss.}
	\label{fig:freedom}
\end{figure}

The operational workflow of FREEDOM is detailed below:
\begin{enumerate}
    \item \textbf{Decoupled CF Branch}: Learnable user ID embeddings $\mathbf{h}_u$ and item ID embeddings $\mathbf{h}_i$ are initialized via Xavier initialization. The collaborative filtering representations (interactive view) are learned by propagating these embeddings through a 4-layer LightGCN over the user-item interaction graph with edge dropout, followed by mean pooling to produce final representations $\mathbf{e}_{u, \text{CF}}$ and $\mathbf{e}_{i, \text{CF}}$.
    \item \textbf{Frozen kNN Graph Construction}: Raw image features (CLIP) and text features (TF-IDF) are projected and fused to generate initial modal representations $\mathbf{e}_{i, \text{modal}}^{(0)}$. At initialization, the cosine similarities between these modal vectors are computed. For each item, only the top-$k$ nearest neighbors (default $k=10$) are retained, and diagonal self-loops are set to zero. The resulting graph is symmetrized and normalized as $\mathbf{A}^{knn} = \mathbf{D}_{knn}^{-1/2} \mathbf{W}^{knn} \mathbf{D}_{knn}^{-1/2}$. This graph structure is frozen (i.e., registers no gradients) during the training phase, serving as a denoised structural prior.
    \item \textbf{Content Propagation}: Item content representations (semantic view) are computed by propagating the projected modal embeddings over the frozen kNN graph $\mathcal{G}_{knn}$ for 4 layers:
    \begin{equation}
        \mathbf{e}_{i, \text{content}}^{(k+1)} = \sum_{j \in \mathcal{N}_k(i)} A^{knn}_{ij} \mathbf{e}_{j, \text{content}}^{(k)},
    \end{equation}
    where $\mathbf{e}_{i, \text{content}}^{(0)} = \mathbf{e}_{i, \text{modal}}$. The final semantic representations $\mathbf{e}_{i, \text{content}}$ are obtained via mean pooling across all layers.
    \item \textbf{InfoNCE Contrastive Alignment}: To align the collaborative interactive view and the semantic content view, FREEDOM applies an InfoNCE loss over positive item pairs (the same item under different views) and negative pairs (different items within the mini-batch $\mathcal{B}$):
    \begin{equation}
        \mathcal{L}_{CL} = -\sum_{i \in \mathcal{B}} \ln \frac{\exp\left(s(\mathbf{e}_{i, \text{CF}}, \mathbf{e}_{i, \text{content}}) / \tau\right)}{\sum_{j \in \mathcal{B}} \exp\left(s(\mathbf{e}_{i, \text{CF}}, \mathbf{e}_{j, \text{content}}) / \tau\right)},
    \end{equation}
    where $s(\cdot, \cdot)$ is the cosine similarity function, and $\tau = 0.2$ is the temperature hyperparameter.
    \item \textbf{Model Fusion and Training Loss}: The final item representation is the sum of the interactive and semantic views:
    \begin{equation}
        \mathbf{e}_i = \mathbf{e}_{i, \text{CF}} + \mathbf{e}_{i, \text{content}}.
    \end{equation}
    The overall model parameters are optimized via:
    \begin{equation}
        \mathcal{L}_{total} = \mathcal{L}_{BPR} + \lambda_1 \mathcal{L}_{reg} + \lambda_2 \mathcal{L}_{CL},
    \end{equation}
    where the contrastive weight $\lambda_2$ is set to $0.1$.
\end{enumerate}

